\chapter{Robustness}
\label{ch:robustness}

The headline result in \cref{ch:results} is that High-AI firms exhibit
a negative differential cumulative-return response to contractionary
\texttt{path\_3f} surprises that becomes statistically distinguishable
from zero at $h \approx 30$ and grows through $h = 60$, reaching
$-7.71$~pp at $h = 60$ in the preferred specification. This chapter
asks how much of that conclusion survives reasonable alternatives to
the specification choices that produced it. \Cref{sec:rob-altshocks}
replaces the forward-guidance factor with the target-rate and
balance-sheet factors; \cref{sec:rob-pathinfo} decomposes the path
factor into monetary and information components;
\cref{sec:rob-altai} swaps the strict$+$ML AI classification for the
original broad keyword grouping; \cref{sec:rob-altoutcomes} replaces
the raw cumulative return with two market-adjusted abnormal returns;
\cref{sec:rob-samplewindows} re-estimates the preferred specification
on nine alternative sample windows; and \cref{sec:rob-inference}
checks the inference under four alternative variance estimators and
a wild-cluster bootstrap. Two complementary checks --- a
second-moment outcome and a sign-decomposition of the shock --- are
reported in \cref{app:rob-volatility,app:rob-asymmetry} because they
speak to mechanism rather than to the existence of the headline
differential. Throughout, the outcome, fixed effects, clustering,
scaling, and active sample are held fixed at the preferred-spec
values unless explicitly varied.

\section{Alternative shocks}
\label{sec:rob-altshocks}

The preferred specification uses \texttt{path\_3f}, the
forward-guidance factor of \citet{swanson2021}. A natural concern is
whether the within-sector AI differential identified in
\cref{ch:results} reflects a generic equity response to any FOMC
announcement surprise, in which case \texttt{target\_3f}
(current-meeting target-rate surprise) and \texttt{lsap\_3f}
(large-scale-asset-purchase surprise) should produce comparable AI
differentials. The three factors are nearly orthogonal in the active
event sample (pairwise $|\rho| \leq 0.091$; \cref{tab:shock_corr}),
so differential AI responses are identified from largely independent
variation. This section re-estimates the preferred specification with
each alternative factor in turn, holding outcome, AI grouping, fixed
effects, clustering, and active sample fixed.

\begin{table}[htbp]
\centering\small
\caption[Alternative shocks: key-horizon estimates]{Alternative shocks:
key-horizon estimates}
\label{tab:alt_shocks_key}
\begin{tabular}{lccccccc}\hline\hline
Coefficient & h=0 & h=5 & h=20 & h=30 & h=40 & h=50 & h=60 \\\hline
target\_3f $\times$ High AI &  0.228 &  0.432 &  0.057 &  0.709 & -0.354 & -0.096 & -0.802 \\
 & ( 0.404) & ( 0.678) & ( 0.930) & ( 1.012) & ( 1.433) & ( 1.904) & ( 1.912) \\
target\_3f $\times$ Moderate AI &  0.076 &  0.056 &  0.212 &  1.473 &  0.146 & -0.808 & -1.251 \\
 & ( 0.099) & ( 0.483) & ( 0.820) & ( 0.863) & ( 0.916) & ( 0.738) & ( 1.179) \\
lsap\_3f $\times$ High AI & -0.016 & -0.169 & -0.335 &  0.089 &  1.248 &  1.930 &  1.893 \\
 & ( 0.355) & ( 0.615) & ( 0.626) & ( 0.760) & ( 1.038) & ( 1.068) & ( 1.015) \\
lsap\_3f $\times$ Moderate AI & -0.231 & -0.182 &  0.569 &  1.030 &  0.941 &  1.192 &  1.488 \\
 & ( 0.054) & ( 0.272) & ( 0.562) & ( 0.643) & ( 0.746) & ( 1.132) & ( 1.078) \\
\hline\hline
\multicolumn{8}{l}{\footnotesize Preferred specification: firm FE + Sector $\times$ Event FE; two-way clustered SEs by firm and event in parentheses.}
\multicolumn{8}{l}{\footnotesize Coefficients rescaled to a +25 bp shock: $\hat\beta \times 25 \times 100$, in percentage points.}
\multicolumn{8}{l}{\footnotesize Omitted group: No AI. Active sample: sample\_preferred == 1, 2005--2024, ai\_score\_inherited == 0.}
\end{tabular}

\end{table}

Under \texttt{target\_3f}, neither AI differential is statistically
distinguishable from zero at any key horizon. High AI point estimates
range from $-0.80$ to $+0.86$~pp and switch sign across horizons;
standard errors lie between $0.40$ and $1.91$~pp. The contrast with
\texttt{path\_3f} is sharp: at $h = 60$ the path-shock High AI
differential is $-7.71$~pp (s.e.~$2.64$) against a target-shock
counterpart of $-0.80$~pp (s.e.~$1.91$). The target-shock estimates
are precise enough to rule out a differential of the path-shock
magnitude at the long horizons that drive the headline.

Under \texttt{lsap\_3f}, the AI differentials are also statistically
indistinguishable from zero at every key horizon, but the picture is
less clean. Point estimates are mostly positive (opposite in sign to
the path-shock estimates) and grow to $+1.89$~pp (High AI) and
$+1.49$~pp (Moderate AI) at $h = 60$, with standard errors near
$1.0$~pp. Two caveats apply. First, the \texttt{lsap\_3f} factor is
identified primarily during the zero-lower-bound period
(\cref{sec:emp-shock}). Second, the full-horizon estimates
(\cref{fig:alt_shocks}(b); full table in
\cref{app:additional-results}) reveal isolated medium-horizon ranges
($h \in [42, 47]$) where the $95\%$ confidence interval just excludes
zero on the upside. The lsap-shock evidence is therefore best
described as imprecise rather than tightly null. Crucially, however,
the lsap-shock differentials do not reproduce the negative,
monotonically growing long-horizon profile of the path-shock baseline.

\begin{figure}[htbp]
    \centering
    \begin{subfigure}[t]{0.48\textwidth}
        \centering
        \includegraphics[width=\textwidth]{fig_07_alt_shock_target_path_high_mod_vs_no.pdf}
        \caption{\texttt{target\_3f}}
    \end{subfigure}
    \hfill
    \begin{subfigure}[t]{0.48\textwidth}
        \centering
        \includegraphics[width=\textwidth]{fig_07_alt_shock_lsap_path_high_mod_vs_no.pdf}
        \caption{\texttt{lsap\_3f}}
    \end{subfigure}
    \caption[Alternative shocks: impulse responses]{Alternative
    shocks: differential cumulative-return response to a $+25$ bp
    surprise in the target-rate factor (panel a) and the LSAP factor
    (panel b), by AI exposure group.}
    \label{fig:alt_shocks}
    \floatfoot{\footnotesize Specification as in
    \cref{fig:preferred_path_irf}; shaded bands are pointwise $95\%$
    CIs.}
\end{figure}

The headline result is therefore specific to the forward-guidance
factor rather than a feature of FOMC announcement surprises in
general. The differential response is concentrated in news about the
expected policy-rate path, not current target-rate news or
balance-sheet news.

\section{Monetary versus information component of the path shock}
\label{sec:rob-pathinfo}

The high-frequency identification used in the preferred specification
treats every announcement-window revision in the path factor as a
monetary surprise. \citet{nakamuraSteinsson2018} document that such
revisions may embed central-bank information about the economic
outlook, and \citet{jarocinskiKaradi2020} propose to separate the two
by the sign of the contemporaneous equity reaction: when rates and
equities move in opposite directions the announcement is consistent
with a textbook monetary surprise; when they move together it is
more naturally read as a signal about growth or risk. Following
\cref{sec:lp-data}, the path factor is decomposed event by event into
\texttt{path\_monetary} (path $\times$ SPFUT $< 0$) and
\texttt{path\_information} (path $\times$ SPFUT $> 0$), where SPFUT
is the contemporaneous S\&P~500 futures surprise. Of the $159$
scheduled FOMC events in the active sample, $93$ are classified as
monetary and $63$ as information; three with $\texttt{path\_3f} = 0$
are unclassified. The preferred specification is re-estimated with
four interactions in place of the two, one per (component, AI group)
combination. \citet{bauerSwanson2023} offer an alternative reading of
the same sign restriction; the decomposition is reported here as a
sensitivity check rather than a structural claim.

\begin{table}[htbp]
\centering\small
\caption[Path--information split: key-horizon estimates]{Path--information
split: key-horizon estimates}
\label{tab:path_info_key}
\begin{tabular}{lccccccc}\hline\hline
Coefficient & h=0 & h=5 & h=20 & h=30 & h=40 & h=50 & h=60 \\\hline
path\_monetary $\times$ High AI & -0.171 &  0.083 & -0.048 & -3.322 & -5.655 & -6.198 & -6.662 \\
 & ( 0.635) & ( 0.985) & ( 1.227) & ( 1.841) & ( 2.206) & ( 2.663) & ( 2.993) \\
path\_monetary $\times$ Moderate AI & -0.073 & -0.756 & -1.196 & -2.209 & -1.508 & -2.470 & -2.878 \\
 & ( 0.390) & ( 0.729) & ( 1.061) & ( 1.659) & ( 1.515) & ( 1.593) & ( 1.753) \\
path\_information $\times$ High AI & -0.986 & -1.694 & -2.450 & -7.795 & -6.963 & -11.067 & -12.140 \\
 & ( 0.693) & ( 2.390) & ( 3.728) & ( 4.213) & ( 4.259) & ( 5.761) & ( 5.911) \\
path\_information $\times$ Moderate AI & -0.248 &  2.043 &  2.712 & -0.546 &  2.632 & -0.523 & -0.413 \\
 & ( 0.709) & ( 1.455) & ( 2.590) & ( 3.314) & ( 3.323) & ( 3.146) & ( 3.344) \\
\hline\hline
\multicolumn{8}{l}{\footnotesize Preferred specification: firm FE + Sector $\times$ Event FE; two-way clustered SEs by firm and event in parentheses.}
\multicolumn{8}{l}{\footnotesize Coefficients rescaled to a +25 bp shock: $\hat\beta \times 25 \times 100$, in percentage points.}
\multicolumn{8}{l}{\footnotesize Omitted group: No AI. Active sample: sample\_preferred == 1, 2005--2024, ai\_score\_inherited == 0. Path factor decomposed by sign of contemporaneous SPFUT reaction.}
\end{tabular}

\end{table}

\begin{figure}[htbp]
    \centering
    \begin{subfigure}[t]{0.48\textwidth}
        \centering
        \includegraphics[width=\textwidth]{fig_09_path_monetary_high_mod_vs_no.pdf}
        \caption{Monetary component}
    \end{subfigure}
    \hfill
    \begin{subfigure}[t]{0.48\textwidth}
        \centering
        \includegraphics[width=\textwidth]{fig_09_path_information_high_mod_vs_no.pdf}
        \caption{Information component}
    \end{subfigure}
    \caption[Path--information split: impulse responses]{Path--information
    split: differential cumulative-return response to a $+25$ bp
    surprise in the monetary component (a) and the information
    component (b), by AI exposure group.}
    \label{fig:path_info_split}
    \floatfoot{\footnotesize Specification as in
    \cref{fig:preferred_path_irf}; panels share a common $y$-axis;
    wider information-panel bands reflect the smaller event count
    ($63$ vs.\ $93$).}
\end{figure}

The monetary component reproduces the headline pattern. The High AI
differential under \texttt{path\_monetary} is indistinguishable from
zero through $h = 30$ and turns negative and distinguishable from
zero at $h = 40$ ($-5.66$~pp, s.e.~$2.21$), $h = 50$ ($-6.20$,
s.e.~$2.66$) and $h = 60$ ($-6.66$, s.e.~$2.99$). The corresponding
preferred-spec baseline values are $-5.90$, $-7.15$, and $-7.71$~pp:
the monetary-component estimates lie within the $68\%$ confidence
interval of the baseline at every key horizon. The Moderate AI
differential under \texttt{path\_monetary} remains imprecise and does
not cross conventional thresholds, consistent with the preferred
specification.

The information component is noisier. High AI point estimates are
also negative at every key horizon $h \geq 5$ and grow to $-12.14$~pp
(s.e.~$5.91$) at $h = 60$. Standard errors are roughly two to three
times those of the monetary component because the information shock
is identified off fewer events. Only $h = 60$ has a $95\%$ CI that
cleanly excludes zero ($[-23.7, -0.6]$); the broader picture is one
of same-signed but imprecise estimates rather than a sharp result.
The Moderate AI differential under \texttt{path\_information} is
small and indistinguishable from zero at every key horizon.

The qualitative reading is that the monetary component of the path
surprise reproduces the long-horizon differential response of High AI
firms at magnitudes within the confidence band of the preferred
specification, while the information component delivers same-signed
but imprecise estimates that the data cannot pin down at the key
horizons. The criticism by \citet{bauerSwanson2023}, that what looks
like an information effect can also reflect market learning about the
policy rule, applies and is left unresolved.

\section{Alternative AI classifications}
\label{sec:rob-altai}

The headline result rests on a discrete three-way classification of
firms by pre-ChatGPT AI exposure, constructed in
\cref{sec:data-ai-classification} from a combination of a strict
keyword filter and a machine-learning--assisted refinement. This
section asks whether the long-horizon differential survives the
alternative discrete classification that uses only the original broad
keyword set without the ML pass. A continuous standardised log AI
score is considered in \cref{app:rob-altai-continuous}.

The original broad classification (\texttt{high2}, \texttt{moderate2},
\texttt{no2}) is mechanically more inclusive in the High and Moderate
tiers. Of the $451$ firms in the active sample,
\cref{tab:alt_ai_crosstab} shows that the two classifications agree
for $428$ ($94.9\%$) and disagree for $23$: $16$ strict No-AI firms
are reclassified upward under the broad scheme ($3$ to High, $13$ to
Moderate), $7$ strict Moderate firms are reclassified as broad High,
and no firm is reclassified downward.

\begin{table}[htbp]
\centering\small
\begin{tabular}{lccc|c}\hline\hline
 & Broad: High & Broad: Moderate & Broad: No & Total \\\hline
Strict+ML: High AI &   79 &    0 &    0 &   79 \\
Strict+ML: Moderate AI &    7 &   64 &    0 &   71 \\
Strict+ML: No AI &    3 &   13 &  285 &  301 \\\hline
Total &   89 &   77 &  285 &  451 \\\hline\hline
\multicolumn{5}{l}{\footnotesize Active sample firms (unique firm\_id values). Both classifications computed pre-ChatGPT.}
\end{tabular}

\end{table}

Re-estimating the preferred specification under the broad
classification yields High AI estimates nearly identical to the
strict$+$ML baseline at every key horizon: $-4.14$ vs.\ $-4.47$~pp
at $h = 30$; $-5.90$ vs.\ $-5.79$ at $h = 40$; $-7.71$ vs.\ $-7.75$
at $h = 60$. The $95\%$ confidence intervals overlap completely at
every horizon, and the same set of horizons ($h \geq 30$) exhibits a
differential distinguishable from zero under both classifications.
The Moderate AI estimates differ in level --- the broad scheme
yields smaller coefficients in absolute value, reflecting that some
strict Moderate firms are absorbed into the broad High group --- but
neither produces a Moderate differential distinguishable from zero
at any key horizon. \Cref{tab:alt_ai_key} reports the full set of
four interaction coefficients across the seven key horizons.

\begin{table}[htbp]
\centering\small
\caption[Strict$+$ML versus original broad classification: key-horizon
estimates]{Strict$+$ML versus original broad classification:
key-horizon estimates}
\label{tab:alt_ai_key}
\begin{tabular}{lccccccc}\hline\hline
Coefficient & h=0 & h=5 & h=20 & h=30 & h=40 & h=50 & h=60 \\\hline
Strict+ML $\times$ High AI & -0.323 & -0.224 & -0.467 & -4.139 & -5.898 & -7.153 & -7.711 \\
 & ( 0.529) & ( 0.872) & ( 1.179) & ( 1.747) & ( 2.067) & ( 2.410) & ( 2.644) \\
Strict+ML $\times$ Moderate AI & -0.106 & -0.242 & -0.469 & -1.865 & -0.706 & -2.113 & -2.420 \\
 & ( 0.323) & ( 0.579) & ( 0.887) & ( 1.468) & ( 1.351) & ( 1.337) & ( 1.584) \\
Original broad $\times$ High AI & -0.347 & -0.738 & -0.870 & -4.465 & -5.790 & -7.382 & -7.752 \\
 & ( 0.520) & ( 0.898) & ( 1.233) & ( 1.745) & ( 1.919) & ( 2.134) & ( 2.396) \\
Original broad $\times$ Moderate AI &  0.085 &  0.120 &  0.163 & -0.267 &  0.299 & -0.177 & -0.656 \\
 & ( 0.273) & ( 0.524) & ( 0.718) & ( 1.263) & ( 1.264) & ( 1.329) & ( 1.220) \\
\hline\hline
\multicolumn{8}{l}{\footnotesize Preferred specification: firm FE + Sector $\times$ Event FE; two-way clustered SEs by firm and event in parentheses.}
\multicolumn{8}{l}{\footnotesize Coefficients rescaled to a +25 bp shock: $\hat\beta \times 25 \times 100$, in percentage points.}
\multicolumn{8}{l}{\footnotesize Omitted group: No AI (no\_ai for Strict+ML; no2 for Original broad). Active sample: sample\_preferred == 1, 2005--2024, ai\_score\_inherited == 0.}
\end{tabular}

\end{table}

The $94.9\%$ firm-level agreement limits what this comparison shows:
the marginal firms reclassified by the ML pass do not carry the
result, but the test cannot speak to how the result would change
under a substantively different exposure measure (e.g., a
revenue-based or patent-based measure). The continuous-exposure check
in \cref{app:rob-altai-continuous} yields same-signed estimates of
comparable magnitude, with an SE-collapse caveat under the preferred
two-way clustering at a subset of horizons that is resolved under
firm-only clustering. The headline is therefore not an artifact of
the specific strict$+$ML construction.

\section{Alternative outcomes: cumulative abnormal returns}
\label{sec:rob-altoutcomes}

The headline outcome \texttt{cum\_ret\_h} is unadjusted for market-wide
return variation. A possible concern is that the AI differential
reflects differential loading on aggregate equity shocks rather than
a genuine cross-sectional response to the forward-guidance surprise.
This section re-estimates the preferred specification with two
market-adjusted alternatives: \texttt{car\_mktadj\_h}, firm return
minus the contemporaneous S\&P~500 return, and \texttt{car\_betaadj\_h},
the beta-adjusted CAR using a firm-specific pre-event market beta.
All three regressions are estimated on the intersection of
observations where the three outcomes are jointly non-missing; the
intersection coincides with the headline active sample at impact
($N = 64{,}851$) and shrinks only mechanically with the horizon (to
$N = 63{,}950$ at $h = 60$).

\begin{table}[htbp]
\centering\small
\caption[Alternative outcomes: key-horizon estimates]{Alternative
outcomes: key-horizon estimates}
\label{tab:alt_outcomes_key}
% Bare tabular. Wrap in \begin{table}...\caption{...}\label{tab:alt_outcomes_key}\end{table} at the call site.
\begin{tabular}{lccccccc}\hline\hline
Outcome & h=0 & h=5 & h=20 & h=30 & h=40 & h=50 & h=60 \\\hline
cum\_ret\_h $\times$ High AI & -0.323 & -0.224 & -0.550 & -4.284 & -5.879 & -7.102 & -7.553 \\
 & ( 0.529) & ( 0.872) & ( 1.224) & ( 1.797) & ( 2.137) & ( 2.503) & ( 2.741) \\
cum\_ret\_h $\times$ Moderate AI & -0.106 & -0.242 & -0.447 & -1.746 & -0.744 & -1.745 & -1.933 \\
 & ( 0.323) & ( 0.579) & ( 0.899) & ( 1.517) & ( 1.394) & ( 1.290) & ( 1.565) \\
car\_mktadj\_h $\times$ High AI & -0.323 & -0.224 & -0.546 & -4.281 & -5.875 & -7.093 & -7.543 \\
 & ( 0.529) & ( 0.872) & ( 1.224) & ( 1.797) & --- & --- & ( 2.741) \\
car\_mktadj\_h $\times$ Moderate AI & -0.106 & -0.242 & -0.443 & -1.743 & -0.740 & -1.736 & -1.921 \\
 & ( 0.323) & ( 0.579) & ( 0.899) & ( 1.517) & --- & --- & ( 1.563) \\
car\_betaadj\_h $\times$ High AI & -0.200 & -0.064 & -0.332 & -3.946 & -5.572 & -6.684 & -6.499 \\
 & ( 0.479) & ( 0.815) & --- & ( 1.863) & ( 2.198) & ( 2.420) & ( 2.677) \\
car\_betaadj\_h $\times$ Moderate AI & -0.025 & -0.162 & -0.303 & -1.533 & -0.630 & -1.598 & -1.602 \\
 & ( 0.290) & ( 0.586) & --- & ( 1.544) & ( 1.568) & ( 1.411) & ( 1.652) \\
\hline\hline
\multicolumn{8}{l}{\footnotesize Preferred specification: firm FE + Sector $\times$ Event FE; two-way clustered SEs by firm and event in parentheses.}
\multicolumn{8}{l}{\footnotesize Coefficients rescaled to a +25 bp shock: $\hat\beta \times 25 \times 100$, in percentage points.}
\multicolumn{8}{l}{\footnotesize Omitted group: No AI. Intersection sample of 64851 observations from the active sample where all three outcomes are non-missing.}
\multicolumn{8}{l}{\footnotesize ``---'' denotes horizons where the two-way clustered SE numerically collapses to zero. Firm-clustered SEs at these horizons are well-defined and reported in \cref{tab:car_vce_appendix}.}
\end{tabular}

\end{table}

The High AI estimates are nearly identical across outcomes. At
$h = 60$ the differential is $-7.55$~pp under \texttt{cum\_ret\_h},
$-7.54$ under \texttt{car\_mktadj\_h}, and $-6.50$ under
\texttt{car\_betaadj\_h}; at $h = 40$ the three values are $-5.88$,
$-5.88$, and $-5.57$. The beta-adjusted CAR produces slightly smaller
magnitudes at the longest horizons --- consistent with the
adjustment removing some differential market loading --- but the
long-horizon pattern of a statistically distinguishable negative
differential for High AI firms holds under all three outcomes. A
handful of cells in \cref{tab:alt_outcomes_key} are marked ``---''
where the two-way clustered variance estimator collapses numerically
at scattered horizons under the CAR outcomes; firm-only clustering
yields well-defined standard errors at every horizon and leaves the
substantive picture unchanged
(\cref{tab:car_vce_appendix} in \cref{app:additional-results}). The
headline differential is therefore a cross-sectional response to the
forward-guidance surprise rather than a reflection of differential
market loading.

\section{Sample-window sensitivity}
\label{sec:rob-samplewindows}

The active sample covers $159$ scheduled FOMC announcements between
February 2005 and December 2024. Two concerns motivate a check across
alternative windows. The first is whether the headline result is
driven by the post-ChatGPT period (December 2022 onward), during
which AI-exposed firms experienced a sustained rerating that
coincided with the 2022--2024 tightening cycle; a positive
correlation between AI exposure and rerating in that window could
mechanically generate a negative differential. The second is whether
unusual periods of monetary policy --- the global financial crisis,
the zero-lower-bound era, the COVID-19 episode, and the longer
pre-2005 sample --- materially shift the estimates. The preferred
specification is re-estimated on nine alternative windows defined in
\cref{tab:samplewindow_definitions}; the first is the preferred
sample itself, reported for reference.

\begin{table}[htbp]
\centering\small
\caption[Sample-window definitions]{Sample-window definitions}
\label{tab:samplewindow_definitions}
% Bare tabular. Wrap in \begin{table}...\caption{...}\label{tab:samplewindow_definitions}\end{table} at the call site.
\begin{tabular}{lllrrr}\hline\hline
 Sample & Start & End & Events & Firms & Obs \\\hline
Preferred 2005-2024 & 02 Feb 2005 & 18 Dec 2024 &       159 &       451 &    64,851 \\
Full 1995-2026 (YF-extended) & 22 Aug 1995 & 28 Jan 2026 &       244 &       462 &    92,965 \\
Full 1995-2024 (CRSP only) & 22 Aug 1995 & 18 Dec 2024 &       235 &       462 &    89,149 \\
2005-2026 (YF-extended) & 02 Feb 2005 & 28 Jan 2026 &       168 &       451 &    68,667 \\
2010-2024 & 27 Jan 2010 & 18 Dec 2024 &       119 &       450 &    49,863 \\
2015-2024 & 28 Jan 2015 & 18 Dec 2024 &        79 &       445 &    33,916 \\
Post-ChatGPT (>= 30nov2022, YF-extended) & 14 Dec 2022 & 28 Jan 2026 &        26 &       442 &    11,286 \\
2005 to ChatGPT (< 30nov2022) & 02 Feb 2005 & 02 Nov 2022 &       142 &       448 &    57,381 \\
Preferred excl. 2020 & 02 Feb 2005 & 18 Dec 2024 &       152 &       451 &    61,825 \\
\hline\hline
\multicolumn{6}{l}{\footnotesize Active sample gate: sample\_preferred == 1, ai\_score\_inherited == 0, and the date filter implied by the window.}
\multicolumn{6}{l}{\footnotesize Obs at h = 0 refers to the firm-event panel at the impact horizon; obs counts at longer horizons are mechanically lower as the post-event window required to compute the outcome clips events at the end of the sample.}
\end{tabular}

\end{table}

\begin{table}[htbp]
\centering\small
\caption[Sample-window sensitivity: High AI--No AI key-horizon
estimates]{Sample-window sensitivity: High AI--No AI key-horizon
estimates}
\label{tab:samplewindow_key}
% Bare tabular. Wrap in \begin{table}...\caption{...}\label{tab:samplewindow_key}\end{table} at the call site.
\begin{tabular}{lccccccc}\hline\hline
Sample & h=0 & h=5 & h=20 & h=30 & h=40 & h=50 & h=60 \\\hline
Preferred 2005-2024 & -0.323 & -0.224 & -0.467 & -4.139 & -5.898 & -7.153 & -7.711 \\
 & ( 0.529) & ( 0.872) & ( 1.179) & ( 1.747) & ( 2.067) & ( 2.410) & ( 2.644) \\
Full 1995-2026 (YF-extended) & -0.455 & -1.200 & -1.241 & -4.047 & -4.852 & -5.780 & -6.672 \\
 & ( 0.438) & ( 0.838) & ( 1.355) & ( 1.847) & ( 2.204) & ( 2.391) & ( 2.664) \\
Full 1995-2024 (CRSP only) & -0.503 & -1.043 & -0.586 & -3.550 & -4.508 & -5.261 & -5.974 \\
 & ( 0.446) & ( 0.848) & ( 1.328) & ( 1.824) & ( 2.201) & ( 2.362) & ( 2.650) \\
2005-2026 (YF-extended) & -0.271 & -0.456 & -1.300 & -4.753 & -6.296 & -7.754 & -8.547 \\
 & ( 0.516) & ( 0.884) & ( 1.256) & ( 1.800) & ( 2.110) & ( 2.467) & ( 2.690) \\
2010-2024 & -0.709 & -0.054 &  0.451 & -3.422 & -6.013 & -6.630 & -6.536 \\
 & ( 0.520) & ( 0.885) & ( 1.254) & ( 1.846) & ( 2.026) & ( 2.089) & ( 2.426) \\
2015-2024 & -0.851 &  0.177 &  1.128 & -2.741 & -6.377 & -6.960 & -7.047 \\
 & ( 0.589) & ( 0.970) & ( 1.291) & ( 1.814) & ( 1.968) & ( 1.999) & ( 2.356) \\
Post-ChatGPT (>= 30nov2022, YF-extended) &  0.308 &  1.393 &  0.279 & -3.842 & -7.344 & -9.103 & -10.627 \\
 & ( 0.770) & ( 1.500) & ( 2.099) & ( 2.712) & ( 2.846) & ( 3.149) & ( 2.997) \\
2005 to ChatGPT (< 30nov2022) & -0.631 & -1.855 & -2.705 & -5.554 & -6.297 & -7.601 & -7.974 \\
 & ( 0.661) & ( 1.207) & ( 1.432) & ( 2.022) & ( 2.694) & ( 3.363) & ( 3.887) \\
Preferred excl. 2020 & -0.285 & -0.260 & -0.425 & -4.015 & -5.777 & -7.019 & -7.484 \\
 & ( 0.528) & ( 0.870) & ( 1.167) & ( 1.729) & ( 2.035) & ( 2.391) & ( 2.615) \\
\hline\hline
\multicolumn{8}{l}{\footnotesize Preferred specification: firm FE + Sector $\times$ Event FE; two-way clustered SEs by firm and event in parentheses.}
\multicolumn{8}{l}{\footnotesize Coefficients rescaled to a +25 bp shock: $\hat\beta \times 25 \times 100$, in percentage points. Reports High AI - No AI differential only; Moderate AI estimates in \cref{tab:samplewindow_key_moderate} in \cref{app:additional-results}.}
\multicolumn{8}{l}{\footnotesize Active sample gate identical across windows except for the date filter; sample sizes for each window in \cref{tab:samplewindow_definitions}.}
\end{tabular}

\end{table}

The headline pattern survives in every window tested. At $h = 40$
the High AI--No AI differential ranges from $-4.51$~pp (full
1995--2024) to $-7.34$~pp (post-ChatGPT) and is statistically
distinguishable from zero at the $5\%$ level in every window. At
$h = 60$ the range is from $-5.97$~pp (full 1995--2024) to
$-10.63$~pp (post-ChatGPT), again excluding zero in every window
(\cref{tab:samplewindow_key}). The decisive comparison is the
pre/post-ChatGPT split. The post-ChatGPT point estimate at $h = 60$
is $-10.63$~pp (s.e.~$3.00$) against a pre-ChatGPT estimate of
$-7.97$~pp (s.e.~$3.89$); the post estimate is roughly $33\%$ larger
in absolute value but rests on only $26$ events, and the two $95\%$
confidence intervals overlap substantially (post $[-16.5, -4.8]$;
pre $[-15.6, -0.4]$). The pre-ChatGPT subset alone --- $142$ FOMC
events spanning 2005--2022 --- identifies the result without any
2023--2024 observations. The headline is therefore not an artifact
of the post-ChatGPT period. Excluding the COVID year, restricting to
the post-GFC or post-liftoff window, and including or excluding the
2025--2026 Yahoo Finance extension all shift the long-horizon
coefficient by less than $1.7$~pp in absolute value and leave every
window distinguishable from zero. The two pre-2005 extensions
attenuate the result by roughly $1.5$~pp at $h = 60$, most plausibly
because the text-based AI-exposure measure is less informative for
the $1995$--$2004$ period during which AI-relevant firm
characteristics had not yet emerged. The Moderate AI estimates
(reported in \cref{tab:samplewindow_key_moderate} in
\cref{app:additional-results}) are indistinguishable from zero at the
seven key horizons in seven of the nine windows; the two exceptions
are the post-GFC windows at $h = 30$, where the Moderate AI estimate
just crosses the $5\%$ threshold with the same sign as the High AI
estimate.

\Cref{fig:samplewindow_irfs} plots the full impulse responses across
$h = 0,\ldots,60$ under each of the nine windows.

\begin{figure}[htbp]
    \centering
    \begin{subfigure}[t]{0.31\textwidth}
        \centering
        \includegraphics[width=\textwidth]{fig_15_time_sample_preferred_2005_2024.pdf}
        \caption{Preferred 2005--2024}
    \end{subfigure}
    \hfill
    \begin{subfigure}[t]{0.31\textwidth}
        \centering
        \includegraphics[width=\textwidth]{fig_15_time_sample_pre_chatgpt.pdf}
        \caption{Pre-ChatGPT}
    \end{subfigure}
    \hfill
    \begin{subfigure}[t]{0.31\textwidth}
        \centering
        \includegraphics[width=\textwidth]{fig_15_time_sample_post_chatgpt.pdf}
        \caption{Post-ChatGPT}
    \end{subfigure}
    \par\bigskip
    \begin{subfigure}[t]{0.31\textwidth}
        \centering
        \includegraphics[width=\textwidth]{fig_15_time_sample_preferred_no_2020.pdf}
        \caption{Excl.\ 2020}
    \end{subfigure}
    \hfill
    \begin{subfigure}[t]{0.31\textwidth}
        \centering
        \includegraphics[width=\textwidth]{fig_15_time_sample_window_2010_2024.pdf}
        \caption{2010--2024}
    \end{subfigure}
    \hfill
    \begin{subfigure}[t]{0.31\textwidth}
        \centering
        \includegraphics[width=\textwidth]{fig_15_time_sample_window_2015_2024.pdf}
        \caption{2015--2024}
    \end{subfigure}
    \par\bigskip
    \begin{subfigure}[t]{0.31\textwidth}
        \centering
        \includegraphics[width=\textwidth]{fig_15_time_sample_full_1995_2024.pdf}
        \caption{Full 1995--2024}
    \end{subfigure}
    \hfill
    \begin{subfigure}[t]{0.31\textwidth}
        \centering
        \includegraphics[width=\textwidth]{fig_15_time_sample_full_1995_2026.pdf}
        \caption{Full 1995--2026}
    \end{subfigure}
    \hfill
    \begin{subfigure}[t]{0.31\textwidth}
        \centering
        \includegraphics[width=\textwidth]{fig_15_time_sample_extended_2005_2026.pdf}
        \caption{2005--2026}
    \end{subfigure}
    \caption[Sample-window sensitivity: impulse responses across
    nine windows]{Sample-window sensitivity: differential
    cumulative-return response to a $+25$ bp path-shock surprise by
    AI exposure group, at $h = 0,\ldots,60$, under nine alternative
    windows.}
    \label{fig:samplewindow_irfs}
    \floatfoot{\footnotesize Specification as in
    \cref{fig:preferred_path_irf}; $y$-axis ranges differ across
    panels to preserve visual resolution; shaded bands are pointwise
    $95\%$ CIs.}
\end{figure}

\section{Inference sensitivity}
\label{sec:rob-inference}

All estimates in \cref{ch:results} use two-way clustered standard
errors by \texttt{firm\_id} and \texttt{event\_id}. The
event-cluster dimension is the binding small-cluster concern, with
$n = 159$ FOMC events --- enough for asymptotic inference but small
enough to warrant a check against the wild-cluster bootstrap of
\citet{cameronGelbachMiller2008}. The preferred specification is
re-estimated under four alternative variance estimators, with the
bootstrap reported at the four long horizons where the headline
differential is statistically distinguishable from zero.

\begin{table}[htbp]
\centering\small
\caption[Inference sensitivity: High AI differential]{Inference
sensitivity: High AI--No AI differential under alternative variance
estimators}
\label{tab:inference_sens_high}
% Bare tabular. Wrap in \begin{table}...\caption{...}\label{tab:inference_sens_high}\end{table} at the call site.
\begin{tabular}{lccccccc}\hline\hline
Inference scheme & h=0 & h=5 & h=20 & h=30 & h=40 & h=50 & h=60 \\\hline
Coefficient (scaled) & -0.323 & -0.224 & -0.467 & -4.139 & -5.898 & -7.153 & -7.711 \\
\hline
Two-way cluster (firm $\times$ event) [preferred] & ( 0.529) & ( 0.872) & ( 1.179) & ( 1.747) & ( 2.067) & ( 2.410) & ( 2.644) \\
Firm cluster & ( 0.325) & ( 0.670) & ( 1.167) & ( 1.821) & ( 2.061) & ( 2.170) & ( 2.283) \\
Event cluster & ( 0.531) & ( 0.910) & ( 1.167) & ( 1.493) & ( 1.810) & ( 2.254) & ( 2.545) \\
Heteroskedasticity-robust & ( 0.326) & ( 0.718) & ( 1.153) & ( 1.571) & ( 1.794) & ( 1.987) & ( 2.160) \\
\hline
Wild-cluster bootstrap by event: 95\% CI &   &   &   & [-7.24, -1.22] & [-9.43, -2.36] & [-11.56, -2.75] & [-12.68, -2.74] \\
\hline\hline
\multicolumn{8}{l}{\footnotesize Point estimate $\hat\beta \times 25 \times 100$ (in percentage points) is constant across rows by construction; only the variance estimator differs.}
\multicolumn{8}{l}{\footnotesize Two-way clustering by firm and event is the preferred specification. Firm- and event-only clustering are nested in the two-way scheme. Heteroskedasticity-robust is reported as a lower bound on the SE.}
\multicolumn{8}{l}{\footnotesize Wild-cluster bootstrap by event ($9{,}999$ replications, percentile-$t$ CI) addresses the small number of event clusters ($n = 159$). Reported at the four long horizons where the headline differential is statistically distinguishable from zero in the preferred specification.}
\end{tabular}

\end{table}

\begin{figure}[htbp]
    \centering
    \includegraphics[width=\WideFigureWidth]{fig_16_inference_sensitivity_high_ai.pdf}
    \caption[Inference sensitivity: High AI confidence bands]{Inference
    sensitivity: High AI--No AI differential cumulative-return
    response to a $+25$ bp path-shock surprise, with $95\%$ confidence
    bands under each of four variance estimators and discrete
    wild-cluster bootstrap intervals at $h \in \{30, 40, 50, 60\}$.}
    \label{fig:inference_sensitivity}
    \floatfoot{\footnotesize Specification as in
    \cref{fig:preferred_path_irf}; the point estimate is constant
    across schemes; wild-cluster bootstrap by event, $9{,}999$
    replications, percentile-$t$ CI.}
\end{figure}

Three observations stand out. First, the High AI differential at
$h \in \{30, 40, 50, 60\}$ has a $95\%$ CI that excludes zero under
every scheme. Second, the preferred two-way clustered SE is the
\emph{most} conservative of the four standard estimators at the long
horizons: at $h = 60$, the two-way SE of $2.64$~pp compares to
firm-clustered $2.28$, event-clustered $2.55$, and
heteroskedasticity-robust $2.16$; the same ordering holds at
$h = 40$ and $h = 50$. The preferred specification is therefore not
understating its standard errors relative to standard alternatives.
Third, the wild-cluster bootstrap delivers nearly identical $95\%$
CIs to the two-way asymptotic test: at $h = 60$ the bootstrap CI is
$[-12.68, -2.74]$ with $p = 0.0053$, against an asymptotic two-way
CI of approximately $[-12.89, -2.53]$. The small-cluster correction
does not change the qualitative conclusion at any of the four
horizons where it is reported.

One pattern is worth flagging. At $h = 20$ and $h = 30$ the
event-clustered SE is slightly smaller than the firm-clustered or
robust SE. Clustering normally raises SEs by allowing within-cluster
correlation; the reverse pattern here reflects mild negative
within-event correlation in the residualised returns. The two-way
clustered SE sits between the two single-cluster estimators, as it
should.

\Cref{tab:inference_sens_moderate} in \cref{app:additional-results}
reports the same comparison for the Moderate AI interaction. The
Moderate AI point estimate is indistinguishable from zero at every
key horizon under every estimator. The wild-cluster bootstrap is not
reported for Moderate AI because the qualitative reading is unchanged
under any standard estimator. The inference underlying the headline
result is therefore not sensitive to the choice of variance
estimator.

\section{Summary and remaining concerns}
\label{sec:rob-summary}

The headline result survives the six checks above. It is specific to
the forward-guidance factor and not to target-rate or balance-sheet
surprises (\cref{sec:rob-altshocks}), concentrated in the monetary
partition of the path factor with same-signed but imprecise estimates
on the information partition (\cref{sec:rob-pathinfo}), nearly
identical under a coarser broad keyword classification
(\cref{sec:rob-altai}), preserved under market-adjusted and
beta-adjusted abnormal returns (\cref{sec:rob-altoutcomes}), present
in all nine alternative sample windows including the pre-ChatGPT
subsample alone (\cref{sec:rob-samplewindows}), and insensitive to
the choice of variance estimator including under a wild-cluster
bootstrap (\cref{sec:rob-inference}). The heterogeneity analysis in
\cref{sec:res-heterogeneity} establishes separately that the
differential is not explained by financial leverage. Complementary
mechanism-oriented evidence in
\cref{app:rob-volatility,app:rob-asymmetry} is qualitatively
consistent with --- but does not establish --- a discount-rate
reading.

Three concerns the chapter does not resolve are worth flagging.
First, the strict$+$ML and original broad classifications agree on
$94.9\%$ of firms, so \cref{sec:rob-altai} cannot speak to how the
result would change under a substantively different AI exposure
measure (e.g., a revenue-based or patent-based measure). Second, the
preferred specification has Sector$\times$Event fixed effects that
partial out industry-event-level variation but no firm-level
time-varying size control; the High AI group includes many of the
largest S\&P~500 firms by market capitalisation in the late sample,
and a size-loading story cannot be ruled out from the checks reported
here. Third, the chapter does not estimate negative-horizon
coefficients to verify that the High AI differential is zero in the
trading days before the announcement. None of these is a known
failure of the headline; each is an extension the data and pipeline
could in principle accommodate.